\section{含参、级数、复数与非初等原函数的完整处理}\label{ic:indefinite_special_complete}

\begin{examtip}这里的“求出”包括给出定义明确、可求导验证的特殊函数表达。对于存在初等原函数的题，先用前述方法约化；特殊函数不是掩盖未完成计算的占位名称。\end{examtip}

\subsection{含参不定积分须固定积分常数}\label{ic:indefinite_special_complete:parameter}\begin{method}{含参不定积分须固定积分常数}\textbf{识别信号} 参数微分能产生对数幂，或消去复杂因子。\par \textbf{操作步骤} 用$F(x,a)=\int_{x_0}^xf(t,a)\dd t$归一化；在紧矩形上验证连续偏导；对参数求导；最后加一个与$x$无关的常数。\par \textbf{疑难点与条件} 未经归一化的原函数可多出$C(a)$，直接对参数求导会漏$C'(a)$。\end{method}

\begin{theorem}{归一化含参原函数}若$f,\partial_af$连续，则$\partial_a\int_{x_0}^xf(t,a)\dd t=\int_{x_0}^x\partial_af(t,a)\dd t$。对$a$多次求导的版本要求相应偏导连续。\end{theorem}\begin{proof}差商在紧积分区间上一致趋向偏导，故积分与极限可交换。下限$x_0$固定，使所有原函数在$x_0$值为0，不产生未知的参数函数。\end{proof}

\begin{example}{考研风格｜参数微分与退化参数}{难度：$\star\star\star$}\label{ic:ex:indefinite_special_complete:parampower}\textbf{题目} 求$\int x^a\ln x\dd x$，$x>0$，并处理$a=-1$。\par \textbf{思路} 从$\frac{x^{a+1}-1}{a+1}$的参数导数开始。\par \textbf{详细解答} $a\ne-1$时，归一化原函数为$\frac{x^{a+1}\ln x}{a+1}-\frac{x^{a+1}-1}{(a+1)^2}$。若只求普通原函数，可把末尾常数吸收成$\frac{x^{a+1}\ln x}{a+1}-\frac{x^{a+1}}{(a+1)^2}+C$。$a=-1$时令$t=\ln x$得$\frac{(\ln x)^2}{2}+C$，也等于归一化公式的参数极限。\par \textbf{求导验算} 对$x$求导，两项中的$\frac{x^a}{a+1}$相消，余下$x^a\ln x$。\par \textbf{疑难点与方法比较} 先固定常数再对参数求导，退化参数可通过极限或单独换元完成。\end{example}

\begin{example}{竞赛风格｜参数产生$x$因子的积分}{难度：$\star\star\star$}\label{ic:ex:indefinite_special_complete:paramtrig}\textbf{题目} 求$\int x\cos(ax)\dd x$。\par \textbf{思路} 对$\int\sin(ax)\dd x$的归一化式求参数导数。\par \textbf{详细解答} $a\ne0$时$\int_0^x\sin(at)\dd t=\frac{1-\cos ax}{a}$，参数微分得到$\int_0^xt\cos(at)\dd t=\frac{x\sin(ax)}{a}+\frac{\cos ax-1}{a^2}$。故普通原函数$\frac{x\sin(ax)}{a}+\frac{\cos(ax)}{a^2}+C$。$a=0$为$\frac{x^2}{2}+C$。直接分部也得到相同公式。\par \textbf{求导验算} 对$x$求导为$\frac{\sin(ax)}{a}+x\cos(ax)-\frac{\sin(ax)}{a}$。\par \textbf{疑难点与方法比较} 未归一化式中$a^{-2}$的发散可被参数常数抵消，不能据此误判$a=0$的原函数不存在。\end{example}

\subsection{幂级数逐项积分与命名函数}\label{ic:indefinite_special_complete:series}\begin{method}{幂级数逐项积分与命名函数}\textbf{识别信号} 已知幂级数在$|x|<R$收敛。\par \textbf{操作步骤} 先在紧子区间用一致收敛逐项积分；写出原函数的收敛半径；若识别为标准函数，给出定义和归一化；端点另判。\par \textbf{疑难点与条件} 数值拟合的级数不等于精确原函数；不能在收敛圆外直接沿用未延拓的级数。\end{method}

\begin{example}{竞赛风格｜反正切积分函数}{难度：$\star\star\star$}\label{ic:ex:indefinite_special_complete:invlogseries}\textbf{题目} 表达$\int\frac{\arctan x\dd x}{x}$，并给$|x|<1$的级数。\par \textbf{思路} 分子级数除以$x$后逐项积分。\par \textbf{详细解答} 定义\[\operatorname{Ti}_2(x)=\int_0^x\dfrac{\arctan t}{t}\dd t,\qquad \left.\dfrac{\arctan t}{t}\right|_{t=0}:=1.\]由反正切展开，在$|x|<1$有\[\arctan x=\sum_{n\ge0}\dfrac{(-1)^nx^{2n+1}}{2n+1},\qquad \operatorname{Ti}_2(x)=\sum_{n\ge0}\dfrac{(-1)^nx^{2n+1}}{(2n+1)^2}.\]所以答案$\operatorname{Ti}_2(x)+C$；在$x=1$该绝对收敛和是Catalan常数。\par \textbf{求导验算} 级数在紧子区间逐项求导恢复$\frac{\arctan x}{x}$，实轴上直接由积分定义验证。\par \textbf{疑难点与方法比较} 本题用精确定义完成表达，不强求不存在的简单初等形式。\end{example}

\begin{example}{考研风格｜有对数也可能完全初等}{难度：$\star\star\star$}\label{ic:ex:indefinite_special_complete:logelementary}\textbf{题目} 求$\int\frac{\ln(1-x)\dd x}{x^2}$，$0<x<1$。\par \textbf{思路} 分部后只剩有理积分。\par \textbf{详细解答} 取$u=\ln(1-x),\dd v=x^{-2}\dd x$，得$-\frac{\ln(1-x)}{x}-\int\frac{\dd x}{x(1-x)}$，故答案$-\frac{\ln(1-x)}{x}-\ln x+\ln(1-x)+C$。\par \textbf{求导验算} 导数中的$\frac{1}{x(1-x)}-\frac{1}{x}-\frac{1}{1-x}$抵消。\par \textbf{疑难点与方法比较} 是否需要多重对数取决于分部后出现的核，不取决于题目中是否有对数。\end{example}

\subsection{复指数与实原函数的还原}\label{ic:indefinite_special_complete:complex}\begin{method}{复指数与实原函数的还原}\textbf{识别信号} 指数三角乘积、三角有理式或高次有理分母。\par \textbf{操作步骤} 利用$\e^{ibx}=\cos bx+i\sin bx$合并；积分后取实部虚部；复对数按共轭成对消去虚常数。\par \textbf{疑难点与条件} 复对数分支变化会产生分段常数；只比较导数不能强迫跨奇点常数相同。\end{method}

\begin{example}{考研风格｜多项式乘复指数}{难度：$\star\star\star$}\label{ic:ex:indefinite_special_complete:complexpolynomial}\textbf{题目} 求$\int x\e^x\cos x\dd x$。\par \textbf{思路} 对$\int x\e^{(1+i)x}\dd x$一次分部。\par \textbf{详细解答} 复原函数$\e^{(1+i)x}[\frac{x}{1+i}-\frac{1}{(1+i)^2}]$。用$\frac{1}{1+i}=\frac{1-i}{2}$、$\frac{1}{(1+i)^2}=-\frac{i}{2}$，取实部得$\frac{\e^x}{2}[x\cos x+(x-1)\sin x]+C$。另一法先求$\int\e^x\cos x$再对$x$分部。\par \textbf{求导验算} 求导后正弦项系数为0，余弦项系数为$x$。\par \textbf{疑难点与方法比较} 复方法只是把两条实线性关系合成一条，不改变最终实常数的自由度。\end{example}

\subsection{Ei、li、Si、Ci的识别核}\label{ic:indefinite_special_complete:ei}\begin{method}{Ei、li、Si、Ci的识别核}\textbf{识别信号} 指数、正弦、余弦分别除以自变量；或$\ln x$在分母。\par \textbf{操作步骤} 固定$\operatorname{Si}(x)=\int_0^x\frac{\sin t\,\dd t}{t}$；$\operatorname{Ci}(x)=\gamma+\ln x+\int_0^x\frac{(\cos t-1)\dd t}{t}$（$x>0$）；实$\operatorname{Ei}(x)=\gamma+\ln|x|+\sum_{k\ge1}\frac{x^k}{k\,k!}$（$x\ne0$）。\par \textbf{疑难点与条件} $\operatorname{Ci}'=\frac{\cos x}{x}$，不是$-\frac{\sin x}{x}$；$\int_0^x\frac{\cos t}{t}$不作为普通收敛定义。\end{method}

\begin{example}{竞赛拓展｜指数积分的复合与分部}{难度：$\star\star\star$}\label{ic:ex:indefinite_special_complete:ei_scaling}\textbf{题目} 求$\int\frac{\e^{2x}\dd x}{x}$与$\int\e^x\ln x\dd x$（第二式$x>0$）。\par \textbf{思路} 第一式识别Ei，第二式分部。\par \textbf{详细解答} 第一式$\operatorname{Ei}(2x)+C$，$x$在不含0的区间。第二式取$u=\ln x,\dd v=\e^x\dd x$，得$\e^x\ln x-\operatorname{Ei}(x)+C$。\par \textbf{求导验算} $\operatorname{Ei}'(2x)\cdot2=\frac{\e^{2x}}{x}$；第二式导数中的$\frac{\e^x}{x}$抵消。\par \textbf{疑难点与方法比较} 链式法则的参数2与分母2x相消。\end{example}

\begin{example}{竞赛拓展｜对数积分的反向复合}{难度：$\star\star\star$}\label{ic:ex:indefinite_special_complete:li}\textbf{题目} 求$\int\frac{\dd x}{\ln x}$，$x>0,x\ne1$。\par \textbf{思路} 设$t=\ln x$，分子变成$\e^t$。\par \textbf{详细解答} 积分$\int\frac{\e^t\dd t}{t}=\operatorname{Ei}(\ln x)+C$。可将该实分支记为$\operatorname{li}(x)$，但跨$x=1$时它有对数奇点，两个区间的常数分别选择。\par \textbf{求导验算} 链式求导$\frac{\frac{\e^{\ln x}}{\ln x}\cdot1}{x}=\frac{1}{\ln x}$。\par \textbf{疑难点与方法比较} 不能从0跨过1直接写成普通收敛积分；若用主值必须明确标记。\end{example}

\subsection{Gaussian、erf、erfi与Fresnel}\label{ic:indefinite_special_complete:gauss}\begin{method}{Gaussian、erf、erfi与Fresnel}\textbf{识别信号} 平方指数或二次相位且无内层导数因子。\par \textbf{操作步骤} 缩放成标准核；多项式权重先分部递推，留下基本特殊函数；若分子已有$x$，优先直接平方变量换元。\par \textbf{疑难点与条件} $\operatorname{erfi}$与$\operatorname{erf}$的指数符号相反；Fresnel常见归一化有不同$\frac{\pi}{2}$因子。\end{method}

\begin{example}{竞赛拓展｜Gaussian二阶矩原函数}{难度：$\star\star\star$}\label{ic:ex:indefinite_special_complete:gaussianmoment}\textbf{题目} 求$\int x^2\e^{-x^2}\dd x$。\par \textbf{思路} 用$x\e^{-x^2}\dd x=-\frac{\dd(\e^{-x^2})}{2}$分部。\par \textbf{详细解答} 积分$=-\frac{x\e^{-x^2}}{2}+\frac12\int\e^{-x^2}\dd x=-\frac{x\e^{-x^2}}{2}+\frac{\sqrt\pi\,\operatorname{erf}x}{4}+C$，其中$\operatorname{erf}x=2\pi^{-\frac{1}{2}}\int_0^x\e^{-t^2}\dd t$。\par \textbf{求导验算} 导数$-\frac{\e^{-x^2}}{2}+x^2\e^{-x^2}+\frac{\e^{-x^2}}{2}$。\par \textbf{疑难点与方法比较} 先消掉可分部的多项式权重，再保留基本核；不必每次定义新函数。\end{example}

\begin{example}{竞赛拓展｜Fresnel归一化}{难度：$\star\star\star$}\label{ic:ex:indefinite_special_complete:fresnelscale}\textbf{题目} 求$\int\sin(3x^2)\dd x$。\par \textbf{思路} 将相位缩放为$\frac{\pi t^2}{2}$。\par \textbf{详细解答} 定义$S(z)=\int_0^z\sin(\frac{\pi t^2}{2})\dd t$，则原函数$\sqrt{\frac{\pi}{6}}\,S(\sqrt{\frac{6}{\pi}}\,x)+C$。如果题目改为$\int x\sin(3x^2)\dd x$，则为$-\frac{\cos(3x^2)}{6}+C$。\par \textbf{求导验算} 前式链式求导的两个缩放系数乘积为1，后式由内层导数$6x$核验。\par \textbf{疑难点与方法比较} 命名函数前必须给出采用的归一化。\end{example}

实变量还采用$\operatorname{erfi}(x)=\frac2{\sqrt\pi}\int_0^x\e^{t^2}\dd t$，所以$a>0$时$\int\e^{ax^2}\dd x=\frac{\sqrt\pi}{2\sqrt a}\operatorname{erfi}(\sqrt a\,x)+C$；负指数相应使用$\operatorname{erf}$。两式均由基本定理与链式法则验证。

\subsection{多重对数的逐阶降阶}\label{ic:indefinite_special_complete:polylog}\begin{method}{多重对数的逐阶降阶}\textbf{识别信号} 出现$\frac{\ln(1-x)}{x}$、$\frac{\operatorname{Li}_s(x)}{x}$或对数乘积。\par \textbf{操作步骤} 在$|x|<1$定义$\operatorname{Li}_s(x)=\sum_{n\ge1}\frac{x^n}{n^s}$，用$\operatorname{Li}_1=-\ln(1-x)$、$\operatorname{Li}_s'=\frac{\operatorname{Li}_{s-1}}{x}$；分部使对数幂下降。\par \textbf{疑难点与条件} 超出单位圆或跨分支切线时需延拓约定；不能把复数恒等式不加条件地当成实式。\end{method}

\begin{example}{竞赛拓展｜三重对数完成对数平方积分}{难度：$\star\star\star$}\label{ic:ex:indefinite_special_complete:polylog3}\textbf{题目} 求$\int\frac{(\ln x)^2\dd x}{1-x}$，$0<x<1$。\par \textbf{思路} 先用$\ln(1-x)$作原函数，再分部两次。\par \textbf{详细解答} 第一次得$-(\ln x)^2\ln(1-x)+2\int\frac{\ln x\ln(1-x)\dd x}{x}$。利用$\frac{\ln(1-x)\dd x}{x}=-\dd\operatorname{Li}_2(x)$，剩余式$=-2\ln x\operatorname{Li}_2(x)+2\operatorname{Li}_3(x)$。总和加$C$即答案。\par \textbf{求导验算} 三项导数中的$\frac{\ln x\ln(1-x)}{x}$和$\frac{\operatorname{Li}_2(x)}{x}$分别成对抵消，余下原式。\par \textbf{疑难点与方法比较} 分部沿$\operatorname{Li}_2\to\operatorname{Li}_3$构成可追踪的层次，不是任意命名。\end{example}

\subsection{Lambert W：用隐式关系换回单元}\label{ic:indefinite_special_complete:lambert}\begin{method}{Lambert W：用隐式关系换回单元}\textbf{识别信号} 出现$W(x)$与$x$的乘除，且能用$x=u\e^u$替换。\par \textbf{操作步骤} 设$u=W(x)$，使用$\dd x=(1+u)\e^u\dd u$；计算后保留所取实分支。\par \textbf{疑难点与条件} 分支点$-\frac{1}{\e}$处$W'$不有限；实主支与下支不能混为同一可微函数。\end{method}

\begin{example}{竞赛拓展｜Lambert W两个基本核}{难度：$\star\star\star$}\label{ic:ex:indefinite_special_complete:lambertpair}\textbf{题目} 求$\int \frac{W(x)\dd x}{x}$与$\int\frac{\dd x}{1+W(x)}$，在相应可微支路。\par \textbf{思路} 统一设$x=u\e^u$。\par \textbf{详细解答} 第一式$\int(1+u)\dd u=u+\frac{u^2}{2}+C=W(x)+\frac{W(x)^2}{2}+C$。第二式$\int\e^u\dd u=\e^{W(x)}+C$；若写成$\frac{x}{W(x)}$，$x=0$处须取极限1。\par \textbf{求导验算} 利用$W'=\frac{W}{x(1+W)}$分别验证；在$W=0$处用连续延拓而非代入$\frac{0}{0}$。\par \textbf{疑难点与方法比较} 隐函数积分与特殊函数积分用同一参数法。\end{example}

\subsection{不完全Beta与Gamma的统一幂权}\label{ic:indefinite_special_complete:incomplete}\begin{method}{不完全Beta与Gamma的统一幂权}\textbf{识别信号} 幂函数乘指数，或$x$与$1+x$的分数幂。\par \textbf{操作步骤} 指数权令$t=ax^n$归到$\gamma(s,z)=\int_0^zt^{s-1}\e^{-t}\dd t$；分式幂令$t=\frac{x}{1+x}$归到$B_z(p,q)$。\par \textbf{疑难点与条件} 先保证定义区间；完全积分的端点参数条件通常强于不完全积分。\end{method}

\begin{example}{竞赛拓展｜指数幂权的完整公式}{难度：$\star\star\star$}\label{ic:ex:indefinite_special_complete:incompletegamma}\textbf{题目} 求$\int x^m\e^{-ax^n}\dd x$，$a,n>0,m>-1,x>0$。\par \textbf{思路} 平方换元的推广$t=ax^n$。\par \textbf{详细解答} 设$s=\frac{m+1}{n}>0$，则原函数$\frac{\gamma(s,ax^n)}{n a^s}+C$。因为$x^m\dd x=\frac{a^{-s}t^{s-1}\dd t}{n}$，完整变换直接得到。若$s$为正整数，$\gamma(s,z)$可用分部递推化为指数乘有限多项式；$s=\frac{1}{2}$化为erf。\par \textbf{求导验算} $\gamma_z(s,z)=z^{s-1}\e^{-z}$，链式乘$anx^{n-1}$后所有幂与常数还原为$x^m\e^{-ax^n}$。\par \textbf{疑难点与方法比较} 这个公式统一许多核，不意味着优先跳过可用的初等换元。\end{example}

这里$B_z(p,q)=\int_0^z t^{p-1}(1-t)^{q-1}\dd t$，$p>0,0<z<1$时允许任意实$q$。例如$x>0,p>0$时，$t=\frac{x}{1+x}$给\[\int x^{p-1}(1+x)^{-p-q}\dd x=B_{\frac{x}{1+x}}(p,q)+C,\]因为$\dd x=\frac{\dd t}{(1-t)^2}$，所有幂化为$t^{p-1}(1-t)^{q-1}$。上限若到1，另要求$q>0$才能得到完全Beta积分。

\subsection{三类椭圆标准化与实参数选择}\label{ic:indefinite_special_complete:elliptic}\begin{method}{三类椭圆标准化与实参数选择}\textbf{识别信号} 根式含两个不同二次因子的乘积，简单代换不能消成有理式。\par \textbf{操作步骤} 先移位缩放把分支点化为标准位置，再$x=a\sin\phi$；分别识别$F,E,\Pi$；退化参数另检。\par \textbf{疑难点与条件} 第一、第二、第三类分别有不同的核，不能全部笼统记作一个椭圆函数。\end{method}

\begin{example}{竞赛拓展｜两个二次因子的三类标准式}{难度：$\star\star\star$}\label{ic:ex:indefinite_special_complete:ellipticthree}\textbf{题目} 设$0<a<b,|x|<a$，将三个积分化为标准形：$I_1=\int\frac{\dd x}{\sqrt{(a^2-x^2)(b^2-x^2)}}$，$I_2=\int\frac{\sqrt{b^2-x^2}\dd x}{\sqrt{a^2-x^2}}$，$I_3$为$I_1$另乘$(1-\frac{\nu x^2}{a^2})^{-1}$。\par \textbf{思路} 令$x=a\sin\phi,k=\frac{a}{b}$。\par \textbf{详细解答} $\phi=\arcsin(\frac{x}{a})$，则$I_1=b^{-1}F(\phi,k)+C$、$I_2=bE(\phi,k)+C$、$I_3=b^{-1}\Pi(\nu;\phi,k)+C$。这里$F_\phi=(1-k^2\sin^2\phi)^{-\frac{1}{2}}$，$E_\phi=\sqrt{1-k^2\sin^2\phi}$，$\Pi_\phi=[(1-\nu\sin^2\phi)\sqrt{1-k^2\sin^2\phi}]^{-1}$。\par \textbf{求导验算} 三式分别乘$\phi'=\frac{1}{\sqrt{a^2-x^2}}$并代$\sqrt{1-k^2\sin^2\phi}=\frac{\sqrt{b^2-x^2}}{b}$，精确还原三核。\par \textbf{疑难点与方法比较} 第三类需避开$1-\frac{\nu x^2}{a^2}=0$；$k=0,1$的退化不能无条件当作一般情形。\end{example}

\begin{example}{竞赛拓展｜四次根式化实模数}{难度：$\star\star\star$}\label{ic:ex:indefinite_special_complete:quarticelliptic}\textbf{题目} 表达$\int\frac{\dd x}{\sqrt{1-x^4}}$，$|x|<1$。\par \textbf{思路} 选择$\sin\phi=\frac{\sqrt2x}{\sqrt{1+x^2}}$使模数为实数。\par \textbf{详细解答} 取$\phi=\arcsin(\frac{\sqrt2x}{\sqrt{1+x^2}})$，则原函数$\frac1{\sqrt2}F(\phi,\frac{1}{\sqrt2})+C$。这避免将$1-x^4=(1-x^2)(1+x^2)$直接记成虚模数而未作说明。\par \textbf{求导验算} $\phi'=\frac{\sqrt2}{(1+x^2)\sqrt{1-x^2}}$，$\sqrt{1-\frac12\sin^2\phi}=\frac{1}{\sqrt{1+x^2}}$，相乘得$\frac{1}{\sqrt{1-x^4}}$。\par \textbf{疑难点与方法比较} 本题保留真正椭圆积分；与分子匹配导数的伪椭圆题分开识别。\end{example}

\subsection{初等表达的边界与最终验算}\label{ic:indefinite_special_complete:boundary}\begin{method}{初等表达的边界与最终验算}\textbf{识别信号} 已完成有理化、分部与代数约化，但仍剩标准非初等核。\par \textbf{操作步骤} 给出已定义的特殊函数或确定的归一化积分；核验导数；写出实定义域及分支。计算机的形式输出也按同一标准复核。\par \textbf{疑难点与条件} 不定积分的所有可能技巧不能用有限清单穷举。常用方法可按结构组织为可执行的算法，但并非任意被积函数都有初等原函数。\end{method}

\begin{warning}常见基本非初等核包括$\e^{-x^2}$、$\sin(x^2)$、$\frac{\sin x}{x}$、$\frac{\e^x}{x}$以及一般椭圆微分。表达式复杂不等于非初等，表达式简单也不保证初等。证明某类核不存在初等原函数通常需要Liouville理论等额外工具；“试了几种换元都失败”不能作为不可积证明。\end{warning}
